% from ega4-2-fr.tex
\oldpage[IV-2]{44}
Supposons \emph{a)} vérifiée. Les hypothèses entraînent que
$\mathcal E\otimes_Y\mathcal F$ est un $\mathcal O_X$-Module cohérent
$(0_{\mathrm I},5.3.11\text{ et }5.3.5)$~; ses cycles premiers associés sont
donc les composantes irréductibles de
$\operatorname{Supp}(\mathcal E\otimes_Y\mathcal F)$ (3.1.4) et pour tout
point maximal $x$ de $\operatorname{Supp}(\mathcal E\otimes_Y\mathcal F)$,
$f(x)=y$ est point maximal de $\operatorname{Supp}(\mathcal E)$ et $x$ point
maximal de $\operatorname{Supp}(\mathcal F_y)$ (2.5.5). Comme, en vertu de
(3.3.1) et du fait que la relation
$y\in f(\operatorname{Supp}(\mathcal F))$ entraîne $\mathcal F_y\ne0$
(I, 9.1.13), tout point de
$\operatorname{Ass}(\mathcal E)\cap f(\operatorname{Supp}(\mathcal F))$ est
l'image par $f$ d'un point maximal de
$\operatorname{Supp}(\mathcal E\otimes_Y\mathcal F)$, on voit que la condition
\emph{b)} est vérifiée.

Inversement, supposons \emph{b)} vérifiée, et montrons que si $z$, $z'$ sont
deux points distincts de $\operatorname{Ass}(\mathcal E\otimes_Y\mathcal F)$,
aucun d'eux ne peut être adhérent à l'autre. En premier lieu, si
$f(z)=f(z')=y$, on a $y\in\operatorname{Ass}(\mathcal E)$ et $z$ et $z'$
appartiennent à $\operatorname{Ass}(\mathcal F_y)$ par (3.3.1), d'où
$y\in f(\operatorname{Supp}(\mathcal F))$~; comme par hypothèse aucun des deux
points $z$, $z'$ n'est adhérent à l'autre \emph{dans} $f^{-1}(y)$, aucun d'eux
ne peut être adhérent à l'autre dans $X$. Si $y=f(z)$ et $y'=f(z')$ sont
distincts, ils appartiennent à
$\operatorname{Ass}(\mathcal E)\cap f(\operatorname{Supp}(\mathcal F))$, donc
aucun d'eux ne peut être adhérent à l'autre dans $Y$~; il résulte de la
continuité de $f$ qu'aucun des points $z$, $z'$ ne peut être adhérent à
l'autre dans $X$.

\begin{proposition}[3.3.4]
\label{IV.2.3.3.4-fr}
Soient $X$, $Y$ deux préschémas localement noethériens, $f:X\to Y$ un
morphisme, $\mathcal E$ un $\mathcal O_Y$-Module cohérent, $\mathcal F$ un
$\mathcal O_X$-Module cohérent et $f$-plat. Alors les conditions suivantes
sont équivalentes~:
\begin{enumerate}
\item[\emph{a)}] $\mathcal E\otimes_Y\mathcal F$ est réduit (3.2.2).
\item[\emph{b)}] Pour tout point
$y\in\operatorname{Ass}(\mathcal E)\cap f(\operatorname{Supp}(\mathcal F))$,
$\overline{\{y\}}$ est un cycle premier associé non immergé de $\mathcal E$,
$\operatorname{long}(\mathcal E_y)=1$ et $\mathcal F_y$ est réduit.
\end{enumerate}
\end{proposition}

Supposons \emph{a)} vérifiée. On sait déjà (3.3.3) que pour tout
$y\in\operatorname{Ass}(\mathcal E)\cap f(\operatorname{Supp}(\mathcal F))$,
$\overline{\{y\}}$ est un cycle premier associé non immergé de $\mathcal E$ et
$\mathcal F_y$ est sans cycle premier associé immergé. En outre (2.5.5), pour
tout $x\in\operatorname{Ass}(\mathcal E\otimes_Y\mathcal F)\cap f^{-1}(y)$,
on a
\[
1=\operatorname{long}((\mathcal E\otimes_Y\mathcal F)_x)
=\operatorname{long}(\mathcal E_y)\cdot\operatorname{long}((\mathcal F_y)_x),
\]
donc $\operatorname{long}(\mathcal E_y)
=\operatorname{long}((\mathcal F_y)_x)=1$, ce qui prouve \emph{b)}.

Inversement, supposons \emph{b)} vérifiée~; on sait déjà que tout point
$x\in\operatorname{Ass}(\mathcal E\otimes_Y\mathcal F)$ est point maximal de
$\operatorname{Supp}(\mathcal E\otimes_Y\mathcal F)$, que $y=f(x)$ est point
maximal de $\operatorname{Supp}(\mathcal E)$ et $x$ point maximal de
$\operatorname{Supp}(\mathcal F_y)$ (3.3.1 et 3.3.3)~; en outre il résulte de
l'hypothèse et de (2.5.5) que
$\operatorname{long}((\mathcal E\otimes_Y\mathcal F)_x)=1$, ce qui prouve
\emph{a)}.

\begin{corollary}[3.3.5]
\label{IV.2.3.3.5-fr}
Soient $X$, $Y$ deux préschémas localement noethériens, $f:X\to Y$ un
morphisme plat~; si $Y$ est réduit aux points de $f(X)$ et si $f^{-1}(y)$ est
un $k(y)$-préschéma réduit pour tout $y\in f(X)$, alors $X$ est réduit.
\end{corollary}

Comme le Nilradical $\mathcal N_Y$ est cohérent, l'ensemble des points où $Y$
est réduit est \emph{ouvert} $(0_{\mathrm I},5.2.2)$, et on peut se borner au
cas où $Y$ est réduit. Il suffit alors d'appliquer (3.3.4) à
$\mathcal E=\mathcal O_Y$ et $\mathcal F=\mathcal O_X$.

\begin{proposition}[3.3.6]
\label{IV.2.3.3.6-fr}
Soient $f:X\to S$, $g:Y\to S$ deux morphismes, $\mathcal F$ un
$\mathcal O_X$-Module quasi-cohérent, $\mathcal G$ un $\mathcal O_Y$-Module
quasi-cohérent. On suppose que~: 1\textsuperscript{o} $\mathcal G$ est
$g$-plat~; 2\textsuperscript{o} $X$ est localement noethérien, et pour tout
$s\in S$, $g^{-1}(s)$ est localement noethérien (ce qui aura lieu si $Y$ est
aussi localement noethérien). Soit $Z=X\times_S Y$~; pour tout couple $(x,y)$
tel que $x\in X$, $y\in Y$ et

